\documentclass{article}
\usepackage{graphicx} % Required for inserting images
\usepackage{enumitem}
\usepackage{microtype}
\usepackage{comment}
\usepackage{indentfirst}
\usepackage{amsthm}
\usepackage{amsmath}
\usepackage{amssymb}


\newtheorem{theorem}{Theorem}[section]
\newtheorem{corollary}[theorem]{Corollary}
\newtheorem{lemma}[theorem]{Lemma}
\newtheorem{proposition}[theorem]{Proposition}
\theoremstyle{definition}
\newtheorem{definition}{Definition}[section]
\title{Capstone}
\author{Huxley Lewis}
\date{November 2025}
\begin{document}
\maketitle

 How can the \textbf{meaningful information limits} of generalized \textbf{AI} structures be inferred through the limits of \textbf{epistemology, constraint satisfaction, and information theory?}

\section{Abstract}

AI is an important part of the modern world, and as architectures are changing, and models are becoming increasingly larger, more general, and intelligent, it becomes important to have an intuitive model for the kinds of problems that AI can solve. This paper aims to provide such a framework for dismissing problems as impossible to solve with AI systems that is both workable and intuitive, incorporating a philosophic perspective, as well as a mathematical rigour.

\section{Introduction}
\begin{comment}
ABSTRACT 
// Restate Question as a fact.
// Justify
// // AI has limitations - INTRO
// // How have other people tackled this problem - LIT REVIEW
// // What is the model - METHOD
// // How did it fare? - ANALYSIS
// // How should it continue? - CONCLUSION
INTRO
// AI is a big thing in the world.
// AI's future is so far undetermined.
// There exist fields that can support our understanding of AI's future.
// // Definition of each field.
// // How are these topics used in research currently?
// // How do they apply to this research?
// Give my definition for AI - Structure = Data → Processing → Generation.
LIT REVIEW
// For each of (AI, CSPs, determinism, empiricism) ... 
// // What are the different definitions of important terms?
// // How do the definitions compare?
// // // Similarities / differences

METHOD 
// Draw from the definitions in the LIT REVIEW and the INTRO
// Link the definitions together
// Form a reduction argument for the limits of AI 
// // Therefore it cannot... 
ANALYSIS
// Test the limits - How would they work in reality?
// What are the limitations of this model?
CONCLUSION
// What are the ethical conclusions?
// How should these ideas be taken forward?
\end{comment}

\medskip
Artificial Intelligence, or AI, has grown significantly over the last decade with the rise of large language models and other foundation models \cite{bommasani2021foundation}. These models are examples of neural networks, which learn from data rather than being programmed with exact task-specific behaviours \cite{lecun2015deep}. Scaling studies of language models suggest that performance improves with additional parameters, data, and compute, while also making the cost of training and inference a central practical constraint \cite{kaplan2020scaling,hoffmann2022training}. Although the exact internal mechanisms by which LLMs generalize across new tasks remain only partially understood \cite{bommasani2021foundation}, their shared structure of data, processing, and generation can still be used to frame arguments about the behaviour and limits of AI systems. This structural view of artificial intelligence motivates the use of constraint satisfiability problems, information theory, epistemology, and determinism when reasoning about the meaningful information limitations of generalized AI models.
\smallskip

Epistemology is the study of human knowledge \cite{sep-epistemology}. It has also been used in recent work to analyse AI systems as black boxes and to reason about the limits of explanation, justification, and understanding in machine learning \cite{carabantes2020blackbox, funer2022accuracy}. Because data is a fundamental element of AI systems, epistemic questions about what is known, how it is learned, and what remains inaccessible apply directly to AI. The different methods by which an AI system acquires and represents data also shape its structure, including the contrast between symbolic and connectionist approaches. (citation)
\smallskip

Constraint Satisfiability Problems, or CSPs, are a general framework for representing problems in terms of variables, domains, and constraints \cite{mackworth1977consistency}. This abstraction makes CSPs useful for discussing classes of problems independently of the details of any particular solver \cite{dechter1992backtracking}. CSPs have long been used in symbolic AI to represent relational structure in domain-specific problems while supporting general solving procedures \cite{mackworth1977consistency,dechter1992backtracking}. This style of representation supports mathematical arguments that may extend to AI systems when those systems can be modelled in CSP-like terms.
\smallskip

Information theory is a framework for quantifying information, communication, and uncertainty \cite{shannon1948mathematical}. In machine learning, information-theoretic methods have been used to study how representations preserve, compress, and transform information during learning \cite{tishby2015deep,shwartz2017opening}. Information theory therefore offers one way to analyse the internal behaviour of contemporary connectionist models \cite{shwartz2017opening}.

\section{Background}

\subsection{Defining AI}

\smallskip

Some cite AI as being any algorithm that aims to replicate human's mental faculties, instead of their physical qualities. This definition, some say, (citation) is redundant since all algorithms are automating mental faculties of people. Other definitions aim to avoid redundancy by specifying how algorithms process data as prediction machines (citation). This provides a much better basis for understanding what AI is, some introductory courses to the field will describe AI as general function approximation algorithms (citation). This has been expanded in recent years as AI models have become successful in problems that were not previously understood as functions (citation), such as natural language processing. The ability to gain meaningful information moves beyond pure numeric function approximation has incurred the use of definitions of AGI, or artificial general intelligence, to differentiate simpler AI from more capable AI models. AGI can similarly be defined in many ways, but, as its name implies, is a generalization of AI, where all mental faculties of people are replicated, this was popularized as a concept by Alan Turing, as he describes the Turing test (citation), whereas a human interacts with an algorithm the human cannot differentiate if an output is the result of an algorithm or another human.

These conflicting definitions are now being shaped by current use of AI as the public are maintaining relationships with the field. This motivates the analysis of these definitions for their applicability in supporting modern uses of AI. The aforementioned definition of AI as automation of mental faculties, in this context, gains meaning since it provides implied use cases for AI, as a tool to replace their mental faculties. The definition of AI as a prediction machine can be introduced to a member of the public to provide understanding of the mechanisms by which AI automates mental processes. The final definition of AI pertinent, here, is that of the general function aproximator, where it again implies a certain use case, that it be used to regress or approximate functions, without the condition of knowledge of the structure of the function.

These standard definitions excel at providing intuitive understand for a lay person to understand what AI is, and how it can be used, but they fall short in the ability to capture the limitations of AI. This motivates the position here, that the definition for AI should be structural, instead of useful, and to provide an intuitive construction of this definition.

\subsection{Epistemic Limitations}

Knowledge has been defined as a justified true belief (JTB), although this has been debated by Gettier (citation), as it does not account for cases such as accidental knowledge. This process of knowledge acquisition has many perspectives. Two of these perspectives are that of the empiricists and rationalists. 
Philosophers such as Hume, argue that knowledge is gained through experience, giving us more evidence, thus increasing the likelihood or probability that what we believe is true, justifying it based on inference from our ability to sense reality (citation).
Descartes, alternatively, in his meditations on first philosophy, that the only knowledge that we can have is that which is infallible, constructing his method of doubting to conclude that his ability to doubt is the only infallibility. These can be understood as different premises of the JTB theory, where Descartes argues that justification and truth can only come through rationality, Hume argues that justification and truth come from our experience of reality. These perspectives ask "what can one know?" in the context of their given reality. In AI problems, the reality must be provided, 

As shown in thought experiments such as Newcomb's problem \cite{nozick1969newcomb}, future events are often treated as examples of the unknowable. This is an example of the epistemic limitations of knowledge, where the future is not knowable in the same way as the past or present. This motivates the idea that there are limits to what can be known, and therefore limits to what can be generated by AI systems. The question then becomes, how can we infer these limits from the structure of AI systems and the constraints that govern their processing?

Knowledge, as defined here is not sufficient to describe the use case of AI models, since knowledge is specific to the agent that understands the knowledge. Let meaningful information, then, be defined as a valid collection or arrangement of signs that incur some meaning to a decoder. Where validity is defined as a constraint of the input and output and meaning is assumed\footnotemark to be captured in the signs of the domain, and signs are symbols that incur meaning beyond their literal shape (citation - semiotic theory).
\begin{definition}[Meaningful Information]\label{def:meaningful}
    Meaningful information is a valid assignment of signs.
\end{definition}
\footnotetext{In extensions of this theory, this assumption must be taken into account, that the problem might be solvable on some general AI model.}

\subsection{Defining Terms}\label{terms}

\begin{definition}[Constraint Satisfiability Problem]\label{def:csp}
    A constraint satisfaction problem (CSP) is a tuple $\mathcal{P} = \langle X, D, C, A \rangle$, where:
    \begin{equation}
        X:=\{x_1, \ldots, x_n\}
        \text{ is a set of variables } x_i.
    \end{equation}
    \begin{equation}
    D:=\{D_1, \ldots, D_n\} \text{assigns a domain } D_i \text{ to each variable } x_i.
    \end{equation}
    \begin{equation}
    C:=\{C_1,\ldots,C_m\} \text{ is a set of constraints } c_i.
    \end{equation}
     Each constraint $C_j\ \in C$ is a pair $C_j=\langle S_j,R_j\rangle$, where $S_j \subseteq X$ is the scope and $R_j$ is a set of valid tuples $R_j \subseteq \underset{x_i \in S_j}{\prod}D_i.$
    \begin{equation}
        A:=\left\{\mathbf{x}\in \underset{\ell}{\prod} D_\ell:\forall c,\ \mathbf{x}_{S_c}\in R_c\right\}, |A|>0 \text{ is the set of valid assignments $\mathbf{x}_k$.}
    \end{equation}
    
\end{definition}


\begin{definition}[Probability]
    Let $P_i(a)$ denote the probability $X_i = a$ or $\Pr (X_i=a)$ and be defined as,
\[
P_i(a)=\frac{|\{\mathbf{x} \in A: x_i=a\}|}{|A|},
\] assuming a uniform distribution. \\
    Let $P_{ij}(a,b)$ denote $\Pr (X_i=a,X_j=b)$, and be defined as,
    \[
    P_{ij}(a,b)=\frac{\left|\left\{\mathbf{x}\in A:\ x_i=a,\ x_j=b\right\}\right|}{|A|}.\]
    
\end{definition}
\begin{definition}[Surprisal]
    The surprisal of an outcome $a$ of random variable $X$ is,
    \[
    I_i(a) := -\log_2\!\bigl(P_i(a)\bigr),
    \]
    where $p(x)$ denotes the probability of $a$.
\end{definition}
\begin{definition}[Entropy]
    The entropy of random variable $X$ is,
    \[
    H(X_i) := -\sum_{a \in D_i} P_i(a)\log_2\!\bigl(P_i(a)\bigr).
    \]
    
\end{definition}
\section{The model of AI}

As there are many definitions of AI, discussed above, we could construct many models of problems for AI models. If we assume some given object class of AI problems, we can then define it by predicating its behaviours. Inductively we can construct certain properties of this class, from our definitions of AI, we can find intersectionality to construct the predicates of this class. In all the definitions of AI, there is a common structure of input, processing, and generation. This structure can be used to construct a model for AI that captures the essential elements of AI systems while allowing for a wide range of specific implementations and architectures. This structure of input, processing, and generation is also consistent with the way that many connectionist AI systems are designed and implemented in practice, where data is collected and preprocessed as input, algorithms are applied to process that data, and outputs are generated based on the processed data. Therefore, this model of AI provides a useful framework for reasoning about the capabilities and limitations of AI systems in terms of their ability to generate meaningful information from their inputs.

\subsection{Contextual Implications}
This model has certain immediate implications for the kinds of problems that AI structures can solve. To reference earlier discussion on epistemic limitations, if the problem is not given a context, then, it can be inferred that there is no ability to justify any generation of meaningful information from any given input, since truth and justification are a consequence of the context. This can, by the assumption of a valid arrangement of signs, be inferred to mean that the meaning and context can be captured in the valid arrangement of signs of the domains of both the input and the generation. Therefore, the input and output domains must be definable into a domain of signs, and there must exist a way of arranging those signs in a valid manner. This implies that the input must be meaningful information by definition \ref{def:meaningful}.

\subsection{Relational Implications}
This model also implies that the relationship between the input and the generation is crucial for determining the meaningful information that can be generated by an AI system. Since the ability to process the information from input to generation must exist. If there is no relationship between the input and the generation, then it is impossible for the AI system to process the input. On the other hand, if there is a relationship between the input and the generation, then the AI system can generate meaningful information, but the amount of information that can be generated is limited by the strength of that relationship. If the relationship is too weak, then the AI system is under-constrained, meaning that the entropy at the output is not less than the entropy at the input, and information is not gained. If the relationship is too strong, then the AI system may be over-constrained, meaning that it cannot generate any output at all. Therefore, the meaningful information that an AI system can generate is limited by the balance of constraints that relate the input to the output.

\subsection{Applying the mathematical model of AI}
The results we inferred in the previous sections can be justified on the basis of the definitions outlined in section \ref{terms}. 

This construction can begin with the input as a variable $x$ with a domain $D_x$, and the generation as a variable $y$ with a domain $D_y$. The processing can be represented as a constraint $c$ that relates $x$ and $y$, such that $c = \langle \{x,y\}, R \rangle$, where $R$ is a set of valid pairs $(x,y)$ such that $|R| > 0$ that satisfy the constraint.
\begin{theorem}
    For any constraint system of more than 2 variables there exists another constraint system with only 2 variables such that there exists a bijection between their sets of satisfying assignments.
\end{theorem}
\renewcommand\qed{$\blacksquare$}
\begin{proof}
    Given a general CSP $\mathcal{P} = \langle X, D, C, A \rangle$ with $|X| \geq  2$, we can construct an equivalent CSP $\mathcal{P}' = \langle X', D', C', A' \rangle$ with $|X'| = 2$ as follows:
    \begin{enumerate}
    \item Let $X' := \{x,y\}$ such that $x$ and $y$ are tuples.
    \item Let $x \subset X$, and $y = X-x$.
    \item Let $D' := \{D_x, D_y\}$.
    \item Let $D_x := \underset{v_i \in x}{\prod} D_i$, and $D_y = \underset{v_i \in y}{\prod} D_i$.
    \item let $C' := \{c'\}$ where $c' := \langle (x,y), R' \rangle$ and $R' = \{(a,b) \in D_x \times D_y: a \bigcup b \in A\}$.
    \item \(A':=\left\{\mathbf{x}\in D_x \times D_y:\ \mathbf{x}\in R'\right\}\)
    \end{enumerate}
\end{proof}
\begin{theorem}
For any CSP $\mathcal P=\langle X,D,C,A\rangle$ with $|X|>2$, there exists a CSP
$\mathcal P'=\langle X',D',C',A'\rangle$ with $|X'|=2$ that is equivalent to $\mathcal P$
(up to a bijection of satisfying assignments).
\end{theorem}

\begin{proof}
Choose a nonempty proper subset $U\subset X$, and let $V:=X\setminus U$.
Define two tuple-variables
\[
X' := \{u,v\},\qquad D_u := \prod_{x_i\in U} D_i,\qquad D_v := \prod_{x_i\in V} D_i,
\]
and $D':=\{D_u,D_v\}$.

Define the binary relation
\[
R' := \{(\alpha_U,\alpha_V)\in D_u\times D_v : \alpha\in A\},
\]
where $\alpha_U,\alpha_V$ are restrictions of assignment $\alpha$ to $U,V$.
Let
\[
C' := \{c'\},\quad c'=\langle \{u,v\},R'\rangle,\quad A':=R'.
\]

Now define
\[
\Phi:A\to A',\qquad \Phi(\alpha):=(\alpha_U,\alpha_V).
\]
Because $U\cap V=\varnothing$ and $U\cup V=X$, each $\alpha\in A$ determines exactly one pair
$(\alpha_U,\alpha_V)$, so $\Phi$ is well-defined.

Define inverse map
\[
\Psi:A'\to A,\qquad \Psi(a,b)=a\oplus b,
\]
where $a\oplus b$ is the unique merged assignment on $X$ agreeing with $a$ on $U$ and $b$ on $V$.
By definition of $R'$, every $(a,b)\in A'$ comes from some $\alpha\in A$, hence $a\oplus b\in A$.

Finally, $\Psi(\Phi(\alpha))=\alpha$ for all $\alpha\in A$, and
$\Phi(\Psi(a,b))=(a,b)$ for all $(a,b)\in A'$. So $\Phi$ is a bijection.
Therefore, $\mathcal P$ and $\mathcal P'$ are equivalent via one-to-one correspondence of satisfying assignments.
\end{proof}
\begin{theorem}
For every CSP $\mathcal P=\langle X,D,C,A\rangle$ with $|X|=n\ge 2$,
there exists an equivalent CSP with exactly $2$ variables.
(Equivalence means a bijection between satisfying assignments.)
\end{theorem}

\begin{proof}
We prove by induction on $n=|X|$.

\textbf{Base case ($n=2$).}
The claim is immediate: $\mathcal P$ already has $2$ variables.

\textbf{Inductive hypothesis.}
Assume that every CSP with $n=k$ variables is equivalent to a CSP with $2$ variables.

\textbf{Inductive step ($k+1\to 2$).}
Let $\mathcal P^{k+1}=\langle X^{k+1},D_{k+1},C_{k+1},A_{k+1}\rangle$ with $|X^{k+1}|=k+1$.
Pick distinct variables $u,v\in X^{k+1}$ and create a new tuple-variable
$w=(u,v)$.

Define a new CSP $\mathcal P^k=\langle X^k,D^k,C^k,A^k\rangle$ by:
\[
X^k := (X^{k+1}\setminus\{u,v\})\cup\{w\},\qquad D_w := D_u\times D_v,
\]
and for all other variables keep the same domains.
Define
\[
A^k := \left\{\beta:\exists \alpha\in A^{k+1}\text{ such that }
\beta(w)=(\alpha(u),\alpha(v))\ \text{and}\ 
\beta(t)=\alpha(t)\ \forall t\in X^{k+1}\setminus\{u,v\}\right\}.
\]
(If needed, let $C^k$ be an extensional/global constraint whose valid tuples are exactly $A^k$.)

Now define
\[
\Phi:A^{k+1}\to A^k,\quad
\Phi(\alpha)=\beta
\]
by the rule above (pack $u,v$ into $w$).
Its inverse
\[
\Psi:A^k\to A^{k+1}
\]
unpacks $w=(a,b)$ back to $u=a,\ v=b$, leaving all other coordinates unchanged.
Thus $\Phi$ is a bijection, so $\mathcal P^{k+1}$ and $\mathcal P^k$ are equivalent.

Since $|X^k|=k$, the inductive hypothesis applies to $\mathcal P^k$:
it is equivalent to a $2$-variable CSP. By transitivity of equivalence,
$\mathcal P^{k+1}$ is equivalent to a $2$-variable CSP.

Therefore the claim holds for all $n\ge 2$.
\end{proof}
\begin{theorem}
\[
\forall n\ge2,\ \mathcal P^n=\langle X^n,D^n,C^n,A^n\rangle,\ |X^n|=n
\ \Longrightarrow\ 
\exists\ \widetilde{\mathcal P}^2\ \text{with }|X(\widetilde{\mathcal P}^2)|=2,\ 
\widetilde{\mathcal P}^2\equiv\mathcal P^n.
\]
\end{theorem}

\begin{theorem}
For every CSP $\mathcal P^n=\langle X^n,D^n,C^n,A^n\rangle$ with $|X^n|=n\ge2$,
there exists an equivalent CSP with exactly $2$ variables.
\end{theorem}

\begin{proof}
Induct on $n$.

Base case $n=2$: immediate.

Inductive hypothesis: true for $n=k$.

Inductive step ($k+1$):
Let $\mathcal P^{k+1}=\langle X^{k+1},D^{k+1},C^{k+1},A^{k+1}\rangle$.
Choose distinct $u,v\in X^{k+1}$ and define a tuple-variable $w=(u,v)$.
Set
\[
X^k := (X^{k+1}\setminus\{u,v\})\cup\{w\},\qquad D_w:=D_u\times D_v.
\]
\[
D^k := \{D_w\}\cup\{D_t:t\in X^{k+1}\setminus\{u,v\}\}.
\]
Define
\[
A^k:=\{\beta:\exists\alpha\in A^{k+1},\ \beta_w=(\alpha_u,\alpha_v),
\ \beta_t=\alpha_t,\ \forall t \in X^{k+1}\setminus\{u,v\}\}.
\]

\[
\implies  C^k = \{\langle X^k, A^k\rangle\}
\]

\(\because  P^k := \langle X^k, D^k, C^k, A^k \rangle\).

Define $\Phi:A^{k+1}\to A^k := \Phi(\alpha,
X_i) = \begin{cases}
(\alpha_u,\alpha_v) & X_i=w,\\
\alpha_i & \forall X_i\in X^k\setminus w.
\end{cases}$

Define $\Psi: A^k\to A^{k+1}
\beta_w = (a,b)
:= \Psi(\beta, X_i) = \begin{cases}
(\beta_w) & X_i=w,\\
\beta_i & \forall X_i\in X^k\setminus w.
\end{cases}$.

Then $\Psi\circ\Phi=\mathrm{id}_{A^{k+1}}$ and $\Phi\circ\Psi=\mathrm{id}_{A^k}$, so $A^{k+1}\cong A^k$.

Hence $\mathcal P^{k+1}\equiv\mathcal P^k$. By hypothesis, $\mathcal P^k\equiv$ a 2-variable CSP.
Therefore $\mathcal P^{k+1}\equiv$ a 2-variable CSP.
\end{proof}

\begin{comment}
    

To justify the model of Input $\rightarrow$ Processing $\rightarrow$ Generation, the limitations established previously must be applied to each of these stages. The empirical and rationalist perspectives motivate this paper's definition of AI as a structure of data, processing, and generation. The data corresponds to the input, the processing corresponds to the transformation of that data, and the generation corresponds to the output. This structure captures the essential elements of AI systems while allowing for a wide range of specific implementations and architectures.

To construct a model for AI, we can draw from the definitions of AI, epistemology, and constraint satisfaction problems. The definition of AI as a structure of data, processing, and generation motivates the idea that AI systems can be modelled as CSPs, where the input data corresponds to the variables and domains, the processing corresponds to the constraints, and the generation corresponds to the valid assignments. This allows us to apply the tools of CSPs to reason about the limits of AI systems in terms of their ability to generate meaningful information from their inputs.

There exists a relationship between the input data and the generated output, which can be understood in terms of the constraints that govern the processing of the data. If there are no constraints that relate the input to the output, then it is impossible for the AI system to generate meaningful information. On the other hand, if there are constraints that relate the input to the output, then the AI system can generate meaningful information, but the amount of information that can be generated is limited by the strength of those constraints. If the constraints are too weak, then the AI system is under-constrained, meaning that the entropy at the output is not less than the entropy at the input, and information is not gained. If the constraints are too strong, then the AI system may be over-constrained, meaning that it cannot generate any output at all. Therefore, the meaningful information that an AI system can generate is limited by the balance of constraints that relate the input to the output.


\subsection{Constructing the model for AI}

To justify the model of Input $\rightarrow$ Processing $\rightarrow$ Generation, the limitations established previously must be applied to each of these stages.
\\
The empirical and rationalist perspectives motivate this paper's definition of AI. 

To apply this to AI, we must construct the context or "reality" of our AI. This motivates the idea that we can suffice this context by constructing a domain of both our inputs and our generation.
\begin{enumerate}
    \item Generation from information is impossible if there exists no relationship between the learned data and generated data.
\end{enumerate}
\begin{proposition} Model claims
    \begin{enumerate}
        \item For all AI models, there exists an equivalent constraint satisfaction representation.
        \item For all AI models, the relationship between generation and information is proportionately tractable to the similarity between the information and the generation.
        \item For all AI models, there exists no equivalent process to generating the future to absolute detail at the current speed of time that is tractable.
    \end{enumerate}
\end{proposition}
\begin{proposition}
    Processing requires that a constraint from $x$ to $y$ exists.
    For constraints between a tuple of variables $T$, a single equivalent constraint can be constructed of $T \cup U$ that is equivalent.
    Over constrained, if there exists no generation, for any of the input's domain the problem is invalid.
    Under constrained, if the entropy of the generation is not less than the entropy of the input, the problem is invalid
\end{proposition}

\subsection{The mathematical model for AI}


\begin{proposition}
    Over constrained:
    $
      D_y(x) = \varnothing.
    $
    Under constrained:
    $
      \lvert D_y(x) \rvert = \lvert D_y \rvert .
    $
    Valid constraint effect:
    $
      0 < \lvert D_y(x) \rvert < \lvert D_y \rvert .
    $
    
\end{proposition}
\end{comment}

\section{Conclusion}
\newpage
\bibliographystyle{unsrt} %Used BibTeX style is unsrt
\bibliography{Bibliography}
\end{document}
